\part {Conclusions} {This is not the end. It is not even the beginning
of the end. But it is, perhaps, the end of the beginning. {\bf
Winston Churchill, 1942}}

\chap Conclusions

The result of this thesis has been the production of a portable
\nn\ simulator called fishNET that will operate with networks of any
size, limited only by available memory and processing speed.

Here, we examine some of the factors important in using the program,
namely performance and flexibility.
Also, we recap some uses for fishNET, and the results of parameters
and casualties on network behaviour.
Lastly, some avenues for possible further work are suggested.

\sec Software -- Design and Performance

A large part of the work for this thesis involved writing and
testing a software package called fishNET.
FishNET was written to be a portable, flexible, network simulator,
and was ported to several installations without change.

The design decisions and their end result, the performance of the
product, have been examined in detail.
The large number of floating point calculations required during
learning tends to be the limiting factor involved with using
fishNET, not a lack of memory, especially in PC implementations.

\sec Possible Uses of FishNET

Briefly, the flexibility of \nn s in general was examined by
testing fishNET with a wide range of problems, for which it
successfully generated generalised networks.
The fact that the same program can be quickly ``taught'' to perform
almost any pattern recognition task is an important result.
Hence, \nn s are an important tool for engineers.

\sec Behaviour of \NN s

The behaviour of \nn s with respect to parameter variations and
casualties was investigated.
The problem of dot-matrix encoded English characters was studied in
considerable detail.

The results of these investigations showed that there can be large
variations in learning times, depending upon the learning
parameter~$\varepsilon$ and the number of \neus\ in the middle
layer.
In summary, learning time is shortest for values of~$\varepsilon$ of
between~$0.3$ and~$0.5$.
Also, for these~$\varepsilon$ values, the learning time decreases
slightly as the number of middle layer \neus\ increases.
However, for larger values of~$\varepsilon$ (of between~$0.9$ and $1.0$), the
learning time {\sl increases\/} as the number of \neus\ in the
intermediate layer gets larger, and a very wide range of learning
times was experienced between runs for these values.
This was especially true for~$\varepsilon = 1.0$, which appears to
be on the boundary of convergence for large numbers of middle layer
\neus.

Error rate was also found to be partly dependent upon~$\varepsilon$
and \neus\ in the middle layer, although not as strongly as learning
time.
In fact, for all the experiments with parameter variation, variance
in error rate was only about~20\%.
However, larger values of~$\varepsilon$ showed lower error rates,
and this seems to be a function of the learning time being longer in
these cases.
Longer learning times logically should return smaller errors, and
this was indeed shown to be the case.

Other important discoveries of network behaviour came from
introducing casualties of two types and varying severity into the
network.
First, random damage to connections (weights) had almost no
appreciable impact for casualty rates up to between~20\% and~30\% 
of the total weight population.
Errors increased exponentially beyond~30\% casualties.
This shows that the knowledge of any pattern is distributed amongst
weights in the network.

Second, random removal of middle layer \neus, or the systematic
removal of weights, had more dramatic effects.
When~2 out of~15 middle layer \neus\ were disabled, performance was
still very good.
However, once another \neu\ was removed, performance degraded
rapidly.
Since the number of weights systematically removed in this way was
still much smaller than with random removal, but the effect was
greater, an important conclusion can be drawn.
The knowledge of features of patterns in a \nn\ is localised or
partly localised into certain \neus\ in the middle layer.
Or, expressing this is another way, certain \neus\ in the middle
layer act as feature extraction devices.
This was widely known, but to this author's knowledge had only been
deduced by looking at the structure of weights.
Here it has been done without examining the distribution and
strength of connections, instead by viewing performance with
casualties in the network.

\sec Further Work

As a large part of the time involved with this thesis was required
to write fishNET, there is still a lot of work that could be done
using the developed software tool.
Some ideas for potential enhancements, and further work, are listed
below.

With fishNET itself, a faster processor such as an~80286/7 or~ 80386/7 
with still a reasonably small memory (say under~2~Mbytes)
would make much larger simulations possible.
To save space, all variables of type {\tt double} could be changed
to {\tt float}, to save almost~50\% from memory usage.
This wasn't done because there was no point for this first version, 
where for reasonably small simulations, speed, not space, was the
limiting factor.

Another improvement to fishNET would be a graphical interface, that
could display outputs during learning (although this would slow
things up even more), and connection strengths both during learning
and execution.
This should be reasonably easy, and would probably involve just
changing the output routines used at present to display graphical
instead of numeric data.
(That was the intention of this first version of fishNET.)

More experiments, especially with regard to error rate, are
desirable.
A wider range of parameters and several simulations for each
set of parameters would remove the statistical ``hiccups'' from the
data, similar to the results for learning rate, which are much
better due to averages being used.

Removing \neus\ in the middle layer offers promise as a mechanism
for discovering the structure of feature selection in a
back-propagation \nn.
More work here could lead to insights into which features are extracted,
by examining behaviour with the individual patterns.

Finally, and most difficult, would be a truly parallel processing
version of the back-propagation algorithm, aimed at operation in
an $n$-transputer environment.

\centerline{---{\eightrm THE END} ---}

